\scanpage{62}
D'après b), ces conditions sont équivalentes dans le cas affine. Donc elles sont équivalentes dans le cas général.

\uline{NB} \deleted{Pour pouvoir conclure des \unsure{ctes|ces} conditions qu'il existe $q,\ c$} \inserted{Supposons que $(L,D,b_0)$ appartienne à l'image essentielle} i.e. que \deleted{la condition} l'on ait \deleted{\unsure{$\exists\,D_1$}}
\[D_1=\gamma(b_0)\]
quand peut on conclure qu'il existe $b_1\in\Gamma(S_1,L_1)$ avec $b_1{}_{S_0}=b_0$, $b_1^2=D_1$ ? \deleted{Or les $b_1$ en question} Bien entendu, les $b_1$ en question sont \uline{exactt} les $b_1$ qui relèvent $b_0$ (grâce à $D_1=\gamma(b_0)$), l'obstruction est donc dans $H^1(S_0,L_0)$. On trouve donc

\uline{Corollaire 1} Pour que $(L,D,b_0)$ soit dans l'image essentielle, il \deleted{\unsure{faut et que}} suffit, et il faut lorsque $H^1(S_0,L_0)=0$ \deleted{qu'il} (p.ex. $S_0$ affine, a fortiori si $S$ affine) qu'il existe $b_1\in\Gamma(S_1,L_1)$ tel que
\[b_1^2=D_1,\]

\uline{Corollaire 2} Supposons $S_0$ réduit. Alors le foncteur $\Delta\longmapsto(L,D)$ est pl. fidèle, et son image essentielle est formée des couples $(L,D)$ satisfaisant les \deleted{deux} conditions

\deleted{a) $D_0=D|S_0$ est de la forme $b_0^2$ $(b_0\in\Gamma(S_0,L_0))$ (NB \deleted{\unsure{alors}} $b_0$ est uniquement déterminée \deleted{\unsure{$f$}}, donc son existence est une condition locale sur $S_0$) b) $\gamma(b_0)=D_1$ \deleted{ou encore} \deleted{ou encore} la section \deleted{\unsure{$\gamma(b_0)$}} \deleted{$[\,\unsure{\alpha|a}(b_0,D_1)\in\Gamma(S_0,L_0)$}}
% archive sheet number: 61
